% Experiment report. Fixed sections: see experiments/TEMPLATE/report.tex and experiments/README.md.
% Build: npm run experiment:report -- 15-rl-strategic-vi
\documentclass{../polyexp}
\begin{document}
% generated by scripts/experiment-report.ts from experiment.json and models/<run>/ - do not edit
\providecommand{\ExpHeader}{}\renewcommand{\ExpHeader}{%
{\sffamily\footnotesize\color{expgrey}EXPERIMENT 15 \textperiodcentered\ METHOD \hfill 2026-11-25}\par
{\LARGE\bfseries\sffamily Strategic actions and value iteration on the full integer state}\hfill\verdictbox{expred}{NOT KEPT}\par\smallskip
{\itshape Does a policy from finite-horizon value iteration over the full integer state, choosing one of five strategic actions per turn, score reliably higher than the bot it steers, and higher than the best single action played every turn?}\par\medskip
\begin{tabular}{@{}p{22mm}p{\dimexpr\linewidth-26mm}@{}}
\fact{Change}{Five strategic actions (bot, expand, economy, army, research: move-type weights $\times$2 for one turn); 300,000 games with a random action per turn; model P̂(s'|s,a), r̂; value iteration gives Q, V and a policy; unseen states play the bot.}
\fact{Knobs}{\knob{boost} $\in$ \{2, 4\}, best 2; \knob{minCount} $\in$ \{1, 5, 20, 50\}, best 5}
\fact{Base}{\model{greedy-full-probe}}
\fact{Method}{Value iteration on 300k random-action games; min-count chosen on 2,000 validation games; final paired test on 10,000 evaluation seeds; metric: final raw score, \#171 rule}
\fact{Setups}{Perfection, 16$\times$16 continents, 3 players, Crazy, seat 0 Xin-Xi (human seat), greedy-full-probe in every seat (as experiment 13)}
\fact{Record}{issue \issue{166}; status: done}
\end{tabular}\par\medskip
\begin{center}\begin{tikzpicture}[node distance=5mm]
    \node[flowbox] (n0) {{\sffamily\bfseries\scriptsize\color{expgrey}BASE}\\[1pt]\model{full-probe}};
    \node[flowbox,right=5mm of n0] (n1) {{\sffamily\bfseries\scriptsize\color{expgrey}CHANGE}\\[1pt]\knob{boost}\\{\scriptsize 2, 4}\\\knob{minCount}\\{\scriptsize 1, 5, 20, 50}};
    \node[flowbox,right=5mm of n1] (n2) {{\sffamily\bfseries\scriptsize\color{expgrey}MEASURED}\\[1pt]Paired, 10,000 evaluation seeds\\{\scriptsize 1 setup, metric: final raw score, \#171 rule}};
    \node[flowbox,right=5mm of n2] (n3) {{\sffamily\bfseries\scriptsize\color{expgrey}RESULT}\\[1pt]Mean final score, 10,000 evaluation seeds: bot vs value-iteration policy\\{\scriptsize 7995.0 $\rightarrow$ 8112.0}\\\textbf{+117.0}};
    \node[flowbox,right=5mm of n3] (n4) {{\sffamily\bfseries\scriptsize\color{expgrey}VERDICT}\\[1pt]\textbf{\color{expred}NOT KEPT}};
    \draw[flowarrow] (n0) -- (n1); \draw[flowarrow] (n1) -- (n2); \draw[flowarrow] (n2) -- (n3); \draw[flowarrow] (n3) -- (n4);
\end{tikzpicture}\end{center}
\begin{center}\sffamily\small\begin{tabular}{@{}lrrr@{}}\toprule
Measure & Base & New & Change\\\midrule
Mean final score, 10,000 evaluation seeds: bot vs value-iteration policy & 7995.0 & 8112.0 & +117.0\\
Same seeds: research every turn vs value-iteration policy & 8386.0 & 8112.0 & $-$274.0\\
Seeds better than the bot (of 9,335 decided) & 4368.0 & 4967.0 & +599.0\\
\bottomrule\end{tabular}\end{center}
}
\providecommand{\ExpResults}{}\renewcommand{\ExpResults}{%
{\small\color{expgrey}No runs recorded.}
}

\ExpHeader

\section{Question and hypothesis}
\textbf{The idea (David, 25 Nov 2026, issue \#166).} Use the full integer state of experiment 13 (turn, cities,
income, units, sum of city levels, techs) to play more strategically: learn from play how each strategic choice
changes the state and the score, then run value iteration backwards from turn 30 to get $V$, $Q$ and a policy.

\textbf{What the action check already showed} (before the main run; Setup): of the five actions played every
turn, only \emph{research} changes the outcome clearly: mean final score 8{,}377 against the bot's 7{,}927, better
on 1{,}167 of 2{,}000 seeds and worse on 792. So the bar is not only the bot but also ``research every turn''.

\textbf{Prediction} (written before the main run):
\begin{enumerate}
  \item \emph{The policy beats the bot.} On the 10{,}000 evaluation seeds the value-iteration policy passes the
        \#171 rule against the bot, with a mean gain of 300--600 points (``research every turn'' gains about 450).
  \item \emph{It mostly researches.} The policy picks research on more than 40\% of learner turns; expand, which
        hardly changes the game, matters little.
  \item \emph{Values are optimistic.} The model's $V$(start) at min-count 1 exceeds the policy's real mean final
        score minus the start score by more than 500 points (the maximum over noisy $Q$ estimates is biased
        upwards, and the state is not Markov); the gap shrinks with a higher min-count.
  \item \emph{The state cannot time the actions much better than a fixed rule.} Against ``research every turn''
        on the same seeds, the policy does not pass the \#171 rule (experiment 13: the six integers carry almost
        no information in the first ten turns, where most of the choice is).
\end{enumerate}
\textbf{Against the idea:} if the policy fails against the bot, value iteration on this state does not even find
the one action that is known to help.

\section{Mechanism}
\textbf{Strategic actions.} At the start of each learner turn one action $a \in \{0,\dots,4\}$ is chosen and holds
for the whole turn: it multiplies the bot's move-type weights (\texttt{FullWeights::ty}, the factor on the score of
every move of that type) by the boost 2. 0 \emph{bot}: no change; 1 \emph{expand}: capture, examine;
2 \emph{economy}: build; 3 \emph{army}: train, attack; 4 \emph{research}: research
(\texttt{engine/rlvalue/src/vi.rs}, \texttt{BOOSTED}).

\textbf{Data.} 300{,}000 games in which the learner picks a uniformly random action each turn (fixed by seed and
turn); each row records the state, the score and the action.

\textbf{Value iteration.} For each state $s$ and action $a$ the observed transitions $(r, s')$ form $\hat P(s'\mid
s,a)$ and $\hat r$. Backwards from the last turn: $Q(s,a) = \frac{1}{n(s,a)}\sum \bigl(r + V(s')\bigr)$,
$V(s) = \max_{a:\,n(s,a) \ge m} Q(s,a)$, $V(\text{terminal}) = 0$; the policy is the argmax. A state where no action
has $m$ visits keeps the mean over all its transitions and has no policy entry. When playing, a state without an
entry gets action 0 (the bot).

\textbf{Choice of $m$.} $m \in \{1, 5, 20, 50\}$; each policy plays 2{,}000 validation games (seeds
$700{,}000{,}000\dots$) against the bot on the same seeds; the $m$ with the highest share of better seeds is kept.
Only that one is evaluated on the evaluation seeds.

\section{Setup}
As experiment 13 (David's choice, \#166, 25 Nov): Perfection, 16$\times$16 Continents, 3 players, Crazy, seat 0
Xin-Xi as the human seat, \texttt{greedy-full-probe} in every seat. This is not the decided 3-opponent gate setup,
so \texttt{rl-gate} is not used: \texttt{rl-value eval} plays the policy and the base on the same seeds and applies
the \#171 rule itself (sign test, exact, one-sided 99\%; mean not significantly worse; 10th percentile not lower).
Seeds: training $1\dots300{,}000$; action check $600{,}000{,}000\dots$ (2{,}000); validation
$700{,}000{,}000\dots$ (2{,}000); evaluation $10^9\dots$ (10{,}000).

\textbf{Action check} (each action every turn, 2{,}000 seeds, \texttt{check-boost2.json}, \texttt{check-boost4.json}):

\begin{center}\small
\begin{tabular}{lrrrrr}
\toprule
action (boost 2) & mean score & median & 10th pct & better / worse vs bot \\
\midrule
bot      & 7927 & 8765 & 2634 & -- \\
expand   & 7924 & 8778 & 2634 & 117 / 110 \\
economy  & 7950 & 8800 & 2585 & 941 / 907 \\
army     & 7969 & 8830 & 2805 & 597 / 606 \\
research & 8377 & 9318 & 3214 & 1167 / 792 \\
\bottomrule
\end{tabular}
\end{center}
With boost 4, research gains more (8{,}497) but costs cities (3.25 against 3.57), economy loses (7{,}686), and
expand still changes little; boost 2 is kept.

Sanity checks: unit tests (value iteration picks the better action; min-count fallback; exact sign test against
$P(X\ge 60\mid 100) = 0.02844$; type-7 quantile; random driver reproducible); the bot driver reproduces experiment
13's games row for row.

Commands (from \texttt{engine/}):
\begin{verbatim}
B=./target/release/rl-value
$B collect --driver random --games 300000 --seed-start 1 --out target/rl-value/train-random.tsv
$B fit-vi --train target/rl-value/train-random.tsv --min-count 1,5,20,50 --out target/rl-value/vi
$B eval --policy target/rl-value/vi/policy-mM.tsv --n 2000 --seed-start 700000000 --out ...
$B eval --policy <chosen> --n 10000 --out ../experiments/15-rl-strategic-vi/eval-vs-bot.json
$B eval --policy <chosen> --n 10000 --base-driver fixed:4 --out .../eval-vs-research.json
\end{verbatim}

\section{Results}
\ExpResults

All numbers are copied from \texttt{vi-summary.json}, \texttt{val-m*.json}, \texttt{eval-vs-bot.json} and
\texttt{eval-vs-research.json} (25 Nov 2026). Data: 300{,}000 random-action games, 9{,}044{,}223 rows, 555{,}338
distinct states; collection 1{,}364\,s on 10 threads.

\paragraph{Value iteration and validation} (2{,}000 validation seeds, policy against the bot on the same seeds).
$V$(start) is the model's remaining score at the turn-0 state; the real remaining score is the policy's mean
final score minus the start score of 515.

\begin{center}\small
\begin{tabular}{rrrrrrr}
\toprule
min-count & policy covers & $V$(start) & real remaining & optimism & better / worse & mean, policy vs bot \\
\midrule
1  & 100.0\% & 11{,}745 & 7{,}543 & +4{,}202 & 922 / 923 & 8058 vs 8055 \\
5  & 86.6\%  & 9{,}399  & 7{,}619 & +1{,}780 & 969 / 882 & 8134 vs 8055 \\
20 & 70.6\%  & 8{,}569  & 7{,}554 & +1{,}015 & 917 / 894 & 8069 vs 8055 \\
50 & 56.1\%  & 8{,}156  & 7{,}529 & +627     & 861 / 898 & 8044 vs 8055 \\
\bottomrule
\end{tabular}
\end{center}
``Policy covers'' is the share of training visits in states with a policy entry. Min-count 5 has the highest share
of better seeds (52.4\%) and is the one evaluated.

\paragraph{Final evaluation} (10{,}000 evaluation seeds from $10^9$, \#171 rule computed by \texttt{rl-value eval}):

\begin{center}\small
\begin{tabular}{lrrrrrl}
\toprule
 & mean & median & 10th pct & better / worse / ties & mean $d$ (99\% upper) & rule \\
\midrule
policy (min-count 5)      & 8112 & 9035 & 2740 & & & \\
bot                       & 7995 & 8850 & 2660 & 4967 / 4368 / 665 & $+117.6$ ($+169.8$) & PASS ($p = 3.0\cdot10^{-10}$) \\
research every turn       & 8386 & 9305 & 3230 & 3913 / 5505 / 582 & $-273.5$ ($-233.0$) & FAIL \\
\bottomrule
\end{tabular}
\end{center}
(The ``better / worse'' columns count the seeds where the policy scored higher / lower than that row's player.)

Reading: the value-iteration policy is reliably better than the bot it steers (better on 53.2\% of the decided
seeds), but by only 118 points on average, and it is clearly worse than the fixed rule ``research every turn''
(274 points less, better on only 41.5\%).

\section{Behaviour}
On the evaluation seeds the policy plays bot 29.6\%, expand 16.7\%, economy 12.0\%, army 17.8\% and research 23.9\%
of the learner's turns; 86.7\% of turns are in a state with a policy entry (the rest play the bot). The choice is
close to uniform: value iteration did not single out research, the one action known to help.

\section{Interpretation}
\textbf{Prediction 1 (beats the bot, +300 to +600): half right.} It passes the \#171 rule, but the gain is +118.

\textbf{Prediction 2 (research on more than 40\% of turns): wrong.} 23.9\%, about the uniform 20\%.

\textbf{Prediction 3 (optimistic values, shrinking with min-count): right.} $V$(start) is 4{,}202 points too high at
min-count 1 and 627 at min-count 50.

\textbf{Prediction 4 (not better than research every turn): right}, and more strongly than predicted: the policy is
significantly worse.

\textbf{Why: the noise of $\hat Q$ swamps the effects of the actions.} The remaining score has a standard deviation
of about 2{,}900 points (experiment 13). An action's effect on it is a few hundred points over a whole game, so a
few dozen points per turn. With 300{,}000 games spread over 555{,}338 states and five actions, a typical
$(s,a)$ is seen a handful of times, so the standard error of $\hat Q(s,a)$ is around a thousand points: the
argmax picks whichever action had lucky games, not the better action. That is the maximisation bias the optimism
column shows, and it is why the choice is nearly uniform. The small gain over the bot comes mostly from the turns
where the policy happens to choose research. Exact integers make the state precise but leave each state with too
little data to tell actions apart; the earlier the turn, the less the state says anyway (experiment 13).

\todo{david}{Your reading: agree, disagree, what is missing?}

\section{Threats to validity}
\begin{itemize}
  \item \emph{Not the decided setup}: 2 opponents, all \texttt{greedy-full-probe}, and the opponent model before
        \#174 (the branch starts from \texttt{dev} at 6f99e0e). The \#171 rule was computed by \texttt{rl-value},
        not \texttt{rl-gate}.
  \item \emph{Exploration data}: values are learned from random-action play, which visits other states than the
        policy does; 13.3\% of the policy's turns fall in states without an entry.
  \item \emph{Min-count chosen on 2{,}000 validation games}: the differences between the four values are within
        noise there (52.4\% at best); another choice would change little.
  \item \emph{Expand} hardly differs from the bot and adds a near-duplicate action to the maximum.
\end{itemize}

\section{Verdict and next step}
\textbf{Not kept.} The policy passes the \#171 rule against the bot (on David's setup, not the gate's), but it is
significantly worse than ``research every turn'', a fixed rule with no learning; keeping it would be keeping noise.
Nothing is promoted (it could not be: the setup is not the gate's).

Next steps this result points to:
\begin{itemize}
  \item \emph{Less noise per $(s,a)$}: more games, or fewer and more different actions (drop expand), or share data
        between neighbouring states (a linear $Q$ in the six integers, i.e.\ function approximation instead of a
        table).
  \item \emph{Spending decisions as actions} (asked on \#166): actions that change the state variables directly.
  \item \emph{Research every turn} is a strong, simple baseline that every later policy should be compared with.
\end{itemize}

\section{Notes for the paper}
Claim: value iteration over a tabular model of exact integer features finds a policy that is reliably but only
slightly better than the bot it steers (+118 points, better on 53\% of seeds), and worse than a fixed rule
(``research every turn'', $-274$); the maximisation bias over sparse $\hat Q$ (the model overestimates the start
value by 1{,}780 points at min-count 5) explains why. Carried by the validation table (optimism column) and the
final evaluation table. A good ``what failed and why'' section for the talk.

\todo{david}{Is this worth a place in the talk or report? Where?}

\end{document}
